\documentclass[10pt, conference]{IEEEtran}
\usepackage[utf8]{inputenc}
\usepackage[margin=0.8in]{geometry}
\usepackage{graphicx}
\usepackage{pgffor}
\usepackage{caption}
\usepackage{paracol}
\usepackage{chngcntr}
\usepackage{amsmath}
\usepackage{xcolor}
\counterwithout{figure*}{section} % If using Article
\bibliographystyle{IEEEtran} % Options: alpha, IEEEtran, unsrt, etc.


\newcommand{\singleWidth}{0.6\textwidth}

\captionsetup[figure]{
    font=small,           % Makes text smaller (can use 'footnotesize' for even smaller)
    labelfont=bf,         % Makes "Figure X" bold
    margin=10pt,          % Indents caption from the left/right of the column
    skip=5pt,             % Reduces space between image and caption
    justification=centering % Keeps it neat
}


\title{Literature review on Semantic Segmentation applied to martian terrain navigation datasets}
\author{Group 17: \\ Darlington Nkrumah (202492437) \\
  Greg de Souza (202582225) \\
  Xuan Toan Doan (202583882)}
\date{April 2026}

\begin{document}

\maketitle

\being{abstract}
This report presents a literature review on semantic segmentation tools and it's
application classifying images from the MSL and M2020 martian missions for the purpose
of navigation. The first section explores the definition and basic structure of
some widely used semantic segmentation models, with focus on the DeepLabV3+ model, and it
also presents the difference between pixel-level and region-level loss functions.
The second section explores the effectiveness of such models on performing semantic
segmentation on the AI4Mars and similar independent datasets, presenting the main
challenges to be yet solved in the efforts for a fully autonomous martian rover
in the future.
\end{abstract}

\section{Semantic Segmentation \& Related Models}

\subsection*{Definitions}
The works by Ghosh et Al \cite{2019}, Minaee et Al \cite{Minaee2020},
Sultana et Al \cite{Sultana2020} and Kale et Al \cite{Kale2021 } explore
the evolution and different possible deep learning techniques for image segmentation.

Image segmentation is an image processing technique used to divide an 
image into two or more meaningful regions, or alternatively, the process 
of defining boundaries separating semantic entities in an image. In a 
more technical definition, the task of image segmentation is to assign 
a label to each pixel within an image, with each label connected with 
respect to some property \cite{Ghosh2019}.

Image segmentation can be divided in three categories: Instance 
segmentation, in which different instances of the same semantic label 
are categorized independently; semantic segmentation, which assigns 
different labels for regions with different semantic meaning without 
classifying independent instances independently; and panoptic 
segmentation, which performs both tasks simultaneously \cite{Sultana2020}.

The task of image segmentation has been significantly impacted by the 
development and advancement of deep learning techniques in computer 
vision, in particular after the success of AlexNet in 2012 in using CNN 
for Semantic Segmentation. Since then, numerous Deep Learning models 
for semantic segmentation have been developed with significant 
success \cite{Sultana2020, Kale2021, Minaee2020}.

\subsection*{Convolutional Neural Networks and Residual Networks}

Classification tasks over an image require as output a probability distribution over a (at least)
2D image. Traditional fully connected networks (perceptrons) are unable to capture spatial
relations due to the linearization requirements. Convolutional networks don't suffer from this issue,
since convolutions can be trivially extended over multiple channels it's possible to input images
in the neural network while retaining the spatial relations at all scales. Convolutional networks
use the convolution operation mixed with layers that resize images to try to extract information
on these multiple scales \cite{Ghosh2019}.


Convolutional layers are the basis for a huge class of models for semantic segmentation include 
Convolutional Neural Networks (CNNs), the define boundary between regions by apply kernel-based 
convolutions to capture the positional relations of features in images. 
By adding pooling and unpooling layers that change the original size, 
it allows CNNs to find meaningful information at multiple scales of 
the feature image by following a ``pyramid-like'' design. CNNs are 
among the most successful approaches to semantic 
segmentation \cite{Minaee2020}. Spatial Pyramid Pooling is a pooling
strategy that resizes/pools the image from larger to smaller to
extract information on all scales.


A relevant advance in convolutional neural networks is the idea 
of dilated (or ``atrous'') convolutions \cite{Minaee2020, Chen2017}, 
which add an additional parameter to convolutions that increases the 
perceptive field without adding new parameters,according to the dilation rate $r$,
shown in equation \ref{eq:atrous}, where $x$ represents the input and
$w$ represents the kernel. A visual representation of the effects of the
dilution rate in a $3 \times 3$ kernel is shown in figure \ref{fig:atrous}.
In practice, the dilution rate makes the convolution pick samples farther from
the center of the kernel, ``skipping'' neighbors, this increases the field
of view without increasing the Kernel size.

\begin{equation}
  y_i = \sum_{k=1}^K x[i+rk]w[k]
  \label{eq:atrous}
\end{equation}


\begin{figure}[hb] % [ht] suggests placement "here" or "top"
  \centering
  \includegraphics[width=0.95\linewidth]{DilatedConvolution.png}
  \caption{Visual representation of diluted convolution with a 3x3 Kernel at
    different dilution rates. \cite{Minaee2020}}
  \label{fig:atrous}
\end{figure}

\begin{figure}[hb] % [ht] suggests placement "here" or "top"
  \centering
  \includegraphics[width=0.95\linewidth]{residualBlock.png}
  \caption{Residual Block \cite{He2016}}
  \label{fig:ResBlock}
\end{figure}

A work by He et Al \cite{He2016} in 2016 proposed a new strategy to build deep neural networks
without multiple layers while avoiding the problem of the vanishing gradient, notably,
the degradation caused by increasing the number of layers isn't an over fitting problem. The
paper addresses the gradient degradation by suggesting a Residual Block as part of the
Neural Network structure. In a residual block, the input of an earlier layer is added to the
the output of a deeper layer. So instead of an usual convolutional layer, where the output
of a block of layers can be written as output$ = F(x)$, a residual block has an output of
$F(x) + x$, as shown in figure \ref{fig:ResBlock}. The work by He et al also shows that residual
blocks allow for deeper networks to still improve and learn with less risk of gradient
degradation, training networks with up to 152 layers and still seeing meaningful improvements
in relation to smaller networks.

The basis for the DeepLab family of models are a combination of diluted convolutions for spatial pooling
(instead of non-learning static pooling layers) with residual blocks, to extract/encode the
information of the original input into a multi channel output, and a CNN for decoding
the encoded input and generate the predictions for classification.
This combination of  convolution layers and
residual blocks to extract the features of the original
image in multiple scales and stack them together in a multichannel output. In
traditional pyramid networks, this final layer is already close to the output
as shown in figure \ref{fig:pyramid}. The DeepLab model adds complexity by using
the pyramid-like section of it's design to just encode the information of the
original input, the encoded features are then extracted by another CNN structure
from the multi-channel pyramid output, as represented in figur\ref{fig:deeplab}.

\begin{figure}[hb] % [ht] suggests placement "here" or "top"
  \centering
  \includegraphics[width=0.95\linewidth]{pyramid.png}
  \caption{The PSPN architecture.
    A CNN produces the feature map and a pyramid pooling module aggregates the different
    sub-region representations. Up-sampling and concatenation are used to
    form the final feature representation from which,
    the final pixel-wise prediction is obtained through convolution \cite{Minaee2020}.}
  \label{fig:pyramid}
\end{figure}


\begin{figure}[hb] % [ht] suggests placement "here" or "top"
  \centering
  \includegraphics[width=0.95\linewidth]{deeplab.png}
  \caption{The DeepLabV3+ model. From \cite{Minaee2020}.}
  \label{fig:deeplab}
\end{figure}


\subsection*{Loss Functions}

Another important choice in implementing a deep learning model is the choice 
of the loss function, a metric that evaluates how far a given output is 
from the expected ground truth. A good choice of loss function can 
mitigate usual challenges of image vision tasks, like noisy datasets 
and unbalanced class representation. 
A number of loss functions are presented and evaluated by the work done 
in \cite{Azad2023}. Pixel-level loss functions utilize the information
in all pixels of the prediction to score the model they usually result
in increased accuracy, specially in scenarios where fine-grained accuracy and
recall are necessary.

One of the most common pixel-level loss metrics is the Cross-entropy (CE) loss 
function. It measures the difference between two probability 
distributions given a random variable. In segmentation tasks, this means 
comparing the ground truth with the assigned probability that the pixel
is of that class according to the model (i.e. output of the soft max function).
The formula for the cross entropy is based on Shannon's entropy, and is presented
in a practical formulation in equation \ref{eq:cross_entropy} , $t_n$ is a one-hot
encoded vector that is 1 on the target class and zero on the others, and $\hat{y}$ is the
predicted probability.  In imbalanced datasets, one common approach to minimize the
effects of the imbalance in the training of the model is adding weights that favor the
less represented classes. If $w$ is a vector of the weights of each labeled class,
the weighted cross entropy is presented in equation \ref{eq:weighted_cross_entropy}.


\begin{equation}
  L_{CE} = -\sum_{i=1}^{C} \log(t_i \cdot \hat{y}_i)
  \label{eq:cross_entropy}
\end{equation}


\begin{equation}
  L_{CE} = -\sum_{i=1}^{C} t_i \cdot w \log(t_i \cdot \hat{y}_i)
  \label{eq:weighted_cross_entropy}
\end{equation}

Cross-entropy loss is part of the Pixel level loss functions, this family of loss functions
evaluate each individual pixel as a part of their computation for dissimilarity, and they usually
improve pixel-based assessment. Another approach are region-level loss functions, that focus
on scoring based on the capacity of  a model to adequately separate different semantic regions of
an image. One example of region-level loss function is the Dice Loss function, that scores
similarity based on the intersection between the predicted subset of pixels $Y$ and the ground
truth region $T$, and is written as in equation \ref{eq:dice} for two class classification task, where
$|A|$ represents the size of a set $A$.
Form multi class segmentation tasks, the formula is presented in equation \ref{eq:multiclass_dice},
where $\hat{y}$ represent the prediction probability of each class, and $t$ is the vector that
represents the ground truth. There is a choice in the number of classes considered in the context
of a possible background/negative class, this background class is sometimes excluded from loss computations.

\begin{equation}
  L_{Dice} = 1 - \frac{2 |Y \cap T|}{|Y| + |T|}
  \label{eq:dice}
\end{equation}


\begin{equation}
  L_{Dice} = 1 - \frac{1}{C} \sum_{c=1}^{C} \frac{2 \sum_{i} t_{ic} \hat{y}_{ic}}{\sum_{i} t_{ic} + \sum_{i} \hat{y}_{ic}}
  \label{eq:multiclass_dice}
\end{equation}

A variation of the Dice Loss that offer some advantages in segmentation tasks is the
log cosh Dice Loss, shown in equation \ref{eq:log_dice}, it enhances smoothness and robustness to outliers,
and minimizes the effects of outliers, and it's a better behaved function in terms of derivatives, that helps
the optimization steps. The Dice Loss can also be ``weighted'' in a sense. A generalized version
of the dice loss is defined by the Generalized Wasserstein Dice Loss, the full scope of the
generalization is outside of this work scope, but a less general version of weighted dice loss
was used in one of the training steps, presented in subsection \ref{sec:train}.

\begin{equation}
  L_{log-cosh} = \log(\cosh(L_{Dice})) 
  \label{eq:log_dice}
\end{equation}

\section{Martian Terrain Segmentation Related Work}

The task of using computer vision for autonomous driving capability has 
been explored by NASA since 2004 \cite{Swan2021} with the AutoNav project. 
The advancements in computer vision through deep learning models 
presented a great opportunity to improve the autonomous driving 
capabilities of the two active Martian rovers: Curiosity, launched as 
part of the Mars Science Laboratory (MSL) mission, and Perseverance, 
launched as part of the Mars 2020 (m2020) mission. Since advancements
in computer vision fueled by development of deep-learning techniques,
several research groups have attempted to use pictures from martian images,
both from rovers and satellites, as input in segmentation tasks, with goal of
identifying meaningful geological features or terrain for navigation.
\cite{Swan2021, Atha2022, Goutham2022, Luna2025}.
The AI4Mars dataset is just of the datasets explored by this groups, while
others have independent dataset, or have labeled the available images through
different methods, or even implemented some level of semi-supervised learning
algorithms \cite{Goutham2022}.

The AI4Mars navigational dataset divides martian terrain in four classes \cite{Swan2021} :
\begin{itemize}
\item \textbf{Soil}: Good for navigation
\item \textbf{Bedrock}: Ideal for navigation.  
\item \textbf{Sand}: A less favorable but still viable path choice.
\item \textbf{Big Rock}: Must be avoided while choosing a path.
\end{itemize}

The AI4Mars dataset is a publicly available set of images and semantic 
segmentation masks developed through a citizen science initiative by 
NASA. A number of volunteers were trained to create the masks, and 
their work was reviewed by specialist laborers and condensed into the 
final dataset \cite{Swan2021}. One notable challenge in the dataset is that the Big Rock class is 
underrepresented, consisting of less than 2\% of pixels in the m2020 
dataset, as shown in Figure \ref{fig:abd_pie}.

The work by Swan et al \cite{Swan2021} trained a DeepLabV3+ with ResNet-101 on two versions of the
MSL dataset, created by two  different labeling strategies called \textit{random}, that used one possible
mask from multiple suggested by trained labeler, and the \textit{merged} strategy, where the final
masks was a merged version of all suggestion by trained labelers, compiled by an expert on martian
terrain. Their work showed that the model performed significantly better with the \textit{merged}
masks, with the accuracy for big rocks increasing from 82.83\% to 93.24\%. The final confusion
matrices from this work are presented in table \ref{tab:Swan2021_results}

\begin{table}[h]
  \centering
  \label{tab:msl_confusion_matrices}
  \begin{tabular}{|l|c|c|c|c|}
    \hline
    \textbf{Actual} & \multicolumn{4}{c|}{\textbf{Predicted}} \\ \cline{2-5} 
                    & \textbf{Soil} & \textbf{Bedrock} & \textbf{Sand} & \textbf{Big Rock} \\ \hline
    \hline
    \multicolumn{5}{|c|}{\textit{MSL Random (Overall Accuracy: 94.97\%)}} \\ \hline
    Soil     & 96.00 & 0.31  & 3.69  & 0.00  \\ \hline
    Bedrock  & 6.15  & 90.87 & 2.54  & 0.44  \\ \hline
    Sand     & 0.25  & 3.23  & 96.51 & 0.01  \\ \hline
    Big Rock & 11.67 & 0.03  & 5.48  & 82.83 \\ \hline
    \hline
    \multicolumn{5}{|c|}{\textit{MSL Merged (Overall Accuracy: 96.67\%)}} \\ \hline
    Soil     & 99.10 & 0.32  & 0.57  & 0.01  \\ \hline
    Bedrock  & 3.64  & 94.90 & 0.37  & 1.09  \\ \hline
    Sand     & 0.88  & 5.62  & 93.45 & 0.05  \\ \hline
    Big Rock & 6.76  & 0.00  & 0.00  & 93.24 \\ \hline
  \end{tabular}
  \caption{MSL NAVCAM Confusion Matrix Percentages for \textit{random} and \textit{merged}
  labeling strategies from Swan et al \cite{Swan2021}}
  \label{tab:Swan2021_results}
\end{table}

The work by Atha et Al \cite{Atha2022} explores additional machine learning models for
semantic segmentation of martian terrain in the context navigation, exploring how
different models perform both on the MSL and on the M2020 navigational dataset. Their work
identified two main problems with current models that aim to allow for independent path
making for martian rovers. First, models tent to perform worse as the rovers move, and terrain
around them changes naturally, the current models aren't general enough for the martian surface.
Second, the accuracy necessary for a model to perform independently, that is, without a remote human
operator offering input, is significantly higher than the accuracy obtained by models at the time of
that work.



\begin{table}[h]
  \centering
  \begin{tabular}{|l|l|c|c|c|}
    \hline
    \textbf{Metric} & \textbf{Class} & \textbf{ResNet-101} & \textbf{+ImageNet} & \textbf{+MSL} \\ \hline
    \hline
    Precision & Soil    & 0.933 & 0.999 & 0.916 \\ \cline{2-5}
                    & Bedrock & 0.857 & 0.467 & 0.917 \\ \cline{2-5}
                    & Sand    & 0.664 & 0.667 & 0.754 \\ \cline{2-5}
                    & Bigrock & 0.558 & 0.618 & 0.174 \\ \hline
    \hline
    Recall    & Soil    & 0.746 & 0.596 & 0.838 \\ \cline{2-5}
                    & Bedrock & 0.570 & 0.958 & 0.536 \\ \cline{2-5}
                    & Sand    & 0.975 & 0.952 & 0.957 \\ \cline{2-5}
                    & Bigrock & 0.283 & 0.355 & 0.640 \\ \hline
  \end{tabular}
  \caption{Performance metrics for models trained on the M2020 dataset by Atha et al. \cite{Atha2022}, reorganized for single-column fit.}
  \label{tab:m2020_atha}
\end{table}

For their model, they trained a ResNet-101 from zero with both MSL and M2020 datasets, they also
trained a ResNet-101 from the weights for ImageNet, and they also tried to transfer train the model
trained on the larger MSL and MER datasets to the M2020 dataset. The transfer learning did improve
performance over soil and rock, but decreased accuracy on labeling bedrock. Their results in
relation to the m2020 Dataset are presented in table \ref{tab:m2020_atha}.

Luna et al \cite{Luna2025} implemented ViBEKO's DeepLabV3+ as part
of their efforts to develop FASTNAV, their tool for independent rover navigation. The dataset used
was independently obtained and labeled by them. Their implementation of DeepLabV3+ with ResNet-50
as a backbone achieved 87\% accuracy and 72\% mIoU with a time to prediction of 25ms, relevant
for practical use in a rover.

Another interesting work in an independent dataset is the work of Goutham et all \cite{Goutham2022}
that attempted to apply self-supervised methods to classify martian terrain using transformers on
the AI4Mars dataset. As a first step, they trained a K-means clustering algorithm to explore what
a deep learning model ``would see'' in relation to texture, structure and specific features of the images
in the dataset. They also implemented a convex hull method for similar purposes. In relation to the dataset,
they made significant effort into pre-processing the images and removing any information that wasn't
relevant for the classification task, such as covering the sky and rover in black pixels.
They then trained different versions of the SegFormer based transformer, and obtained mIoU of more than
80\% in all ther mdoels, with accuracy close to 90\%.


\section{Conclusion}

Semantic segmentation is an elaborate and non-trivial task in computer vision, and performing
segmentation on usual terrain such as martian terrain proves a challenge even for state of the
art models. The lack of a very large dataset, and that the planet Mars itself offers a variety
of very diverse terrain also limits the durability of successful models, a highly accurate model
for the current region of a rover might perform significantly worse after the rover traveled enough
for it's immediate terrain to change significantly. This is also reinforced
by the fact that transfer learning from the MSL to the M2020 dataset doesn't seem to be
particularly effective. Fully autonomous navigation would be a great advancement for interplanetary
exploration on Mars, and modern deep learning models have made great stride in providing very
useful navigational tools, but even better and more robust models are necessary for fully
autonomous navigation in martian images.
 
\bibliography{references}  % Use the filename WITHOUT the .bib extension


\end{document}
